Here I mainly follow the discussion in the appendix of
\cite{carrollSpacetimeGeometryIntroduction2014}.

\subsection{Vielbein as non-coordinate Basis}
\subsubsection{Definition of Vielbein}
In GR we know that the coordinate system on the manifold naturally
generates a basis of the tangent space, which is called the coordinate basis.
$ \partial_{\mu}, dx^\mu $
However, we can also choose other basis of the tangent space, we
define a general othornomal basis as:
\defi{
  Othornormal Basis

  we define a set of basis vectors of the tangent space as $ e_a $
  which satisfy:
  \begin{align}
    g(e_a, e_b) = \eta_{ab}
  \end{align}
  where $ \eta_{ab} $ is depend on the spacetime signature.
  Equivantly, we can define a set of orthonormal basis of the
  cotangent space as $ \theta^a $ which satisfy:
  \begin{align}
    \theta^a(e_b) = \delta^a_b
  \end{align}
}
Now we define the Vielbein as:
\begin{itemize}
  \item Components of the coordinate basis in an orthonormal basis: $
    \partial_{\mu} = e_{\mu}{}^a e_a $
  \item Transformation matrix between the coordinate basis and the
    orthonormal basis
\end{itemize}
We can also define the inverse of the vielbein as: $ e^\mu{}_a $ (note
  we always put the coordinate index first and the orthonormal index
second), which satisfy:
\begin{align}
  {e^\mu}_a e_\mu{}^b = \delta_a^b, \quad e_\mu{}^a {e^\nu}_a =
  \delta_\mu^\nu.
\end{align}
We then can see that vielbein and inverse vielbein have the following property:
\begin{enumerate}
  \item    Can be used as basis
    transformation matrices between coordinate basis and orthonormal
    basis in tangent space and cotangent space.
    \begin{align}
      e_a = e^\mu{}_a \partial_{\mu}, \quad \theta^a = e_\mu{}^a
      dx^\mu \quad \partial_{\mu} = e_{\mu}{}^a e_a, \quad dx^\mu =
      e^\mu{}_a \theta^a
    \end{align}
  \item    The orthonormal condition can be written as:
    \begin{align}
      g_{\mu\nu} = e_\mu{}^a e_\nu{}^b \eta_{ab}, \quad \eta_{ab} =
      e^\mu{}_a e^\nu{}_b g_{\mu\nu}
    \end{align}
\end{enumerate}

Now we consider a general tensor field. A tensor can be written in
the coordinate basis or the orthonormal basis or even mixed basis.
For example, a vector field can be written as:
\begin{align}
  V = V^\mu \partial_{\mu} = V^a e_a
\end{align}
A tensor field can be written as:
\begin{align}
  T = T^{\mu\nu} \partial_{\mu} \otimes \partial_{\nu} = T^{ab} e_a
  \otimes e_b = T^{\mu a} \partial_{\mu} \otimes e_a
\end{align}
Moreover we can define a more general raising and lowering index rule:
\begin{itemize}
  \item For a coordinate index, we can raise and lower it with the metric $ g $.
  \item For an orthonormal index, we can raise and lower it with $ \eta $.
\end{itemize}
This rule is exactly consistent with the definition of the vielbein,
since we have:
\begin{align}
  e_\mu{}^a = g_{\mu\nu} e^\nu{}_b \eta^{ab}, \quad e^\mu{}_a =
  g^{\mu\nu} e_\nu{}^b \eta_{ab}
\end{align}
Finally we note that there are two other kinds of understanding of the vielbein:
\begin{itemize}
  \item view $ e_\mu{}^a$ as a set of one forms, which is just the
    basis vectors of the cotangent space $ \theta^a $ so we may
    rename it as $ e^a = e_\mu{}^a dx^\mu $.
  \item view $ e^\mu{}_a $ as a (1,1) tensor field.
\end{itemize}

\subsubsection{Basis Transformations GCT and LLT}

Now we consider basis transformation, for a coordinate basis, the
basis of the tangent space is tight to the coordinate system. Thus,
it changes only when we choose a different coordinate system and the
transformation matrix is given by the Jacobian of the coordinate
transformation. We call this transformation as \textbf{GCT (General
Coordinate Transformation)}.

However, for a orthonormal basis, we have the freedom of doing a
basis transformation without changing the coordinate system. This is because:
\begin{align}
  e^a \to e'^a = \Lambda^a{}_b(x) e^b
\end{align}
is also a valid orthonormal basis as long as $ \Lambda^a{}_b(x) $ is
a local Lorentz transformation, i.e., it satisfies:
\begin{align}
  \Lambda^a{}_c(x) \Lambda^b{}_d(x) \eta_{ab} = \eta_{cd}
\end{align}
This basis transformation of the orthonormal basis is called
\textbf{LLT (Local Lorentz Transformation)}.
note that different from the lorentz transformation in the coordinate basis:
\begin{itemize}
  \item LLT doesn't change the coordinate system, it only changes the
    basis of the tangent space.
  \item LLT is a local transformation, the transformation matrix can
    depend on the spacetime point instead of being a constant.
\end{itemize}
Generally, if we fix the orthonormal basis, then GCT will change the
vielbein yet keep the orthonormal basis unchanged, while LLT is just
a change of the orthonormal basis, thus it will change the vielbein
but keep the coordinate basis unchanged.
\subsection{Covariant Derivative}

\subsubsection{Spin Connection}

Now we consider the covariant derivative of a vector field when we
write it out in the orthormormal basis. To do this, we need to define
a covariant derivative of the orthonormal basis vectors, which lead
us to the spin connection:
\defi{
  Spin Connection

  We define the spin connection $ \omega_\mu{}^a{}_b $ as:
  \begin{align}
    \nabla_\mu e_a = \omega_\mu{}^b{}_a e_b \quad \nabla_\mu e^a
    = -\omega_\mu{}^a{}_b e^b
  \end{align}
}
In fact just as the normal connection, we can derive the covariant
derivative of the cotangent space basis from that of the tangent
space basis, we only need to define the first part.

With this definition, we can write the covariant derivative of a
vector field component in the orthonormal basis as:
\begin{align}
  \nabla_\mu X^a{}_b = \partial_\mu X^a{}_b + \omega_\mu{}^a{}_c
  X^c{}_b - \omega_\mu{}^c{}_b X^a{}_c
\end{align}
We can prove a relation between the spin connection and the normal
connection by evaluatingf the covariant derivative of a general
vector field $ \nabla X $ in both basis and compare the results. We have:
\begin{align}\label{relationconnection}
  \Gamma^\nu_{\mu\lambda} = e^\mu{}_a \partial_\mu e_\lambda{}^a +
  e^\nu{}_a e_\lambda{}^b \omega_\mu{}^a{}_b \quad \omega_\mu{}^a{}_b
  = -e^\nu{}_b \partial_\mu e_\nu{}^a + e^\nu{}_b e_\lambda{}^a
  \Gamma^\lambda_{\mu\nu}
\end{align}
Then we evaluate the covariant derivative of Vielbein, if we view it
as a 1-form, we have:
\begin{align}
  \nabla_\mu e_{\mu}{}^a = \partial_\mu e_\nu{}^a -
  \Gamma^\lambda_{\mu\nu} e_\lambda{}^a
\end{align}
We can plug in the relation between the spin connection and the
normal connection, we have:
\begin{align}
  \nabla_\mu e_{\nu}{}^a = -\omega_\mu{}^a{}_b e_\nu{}^b
\end{align}
We can see that everything is constistent.

\subsubsection{Spin Connection under LLT}

Just as a normal connection spin connection is not a tensor. Yet due
to the definition and linearity of the covariant derivative $ \mu $
index, we can prove that:
\begin{itemize}
  \item Under GCT the spin connection $ \mu $ index transforms as a
    tensor. This means that we can view the spin connection as a one
    form in the coordinate basis with two indexes.
  \item  Yet under LLT the spin connection doesn't transform as a
    tensor, we can prove from the definition or the property of the
    covariant derivative is a tensor under LLT that:
\end{itemize}
\begin{align}
  \omega_\mu{}^{a'}{}_{b'} = \Lambda^{a'}{}_a
  \Lambda^b{}_{b'} \omega_\mu{}^a{}_b + \Lambda^{a'}{}_c \partial_\mu
  \Lambda^c{}_{b'}
\end{align}
we shall notice that $ \Lambda^{a'}{}_a $ and $ \Lambda^b{}_{b'} $
are inverse transformation matrix of each other.

\subsection{Cartan Structure Equations}

Then comes the most important part of the vielbein formalism, which
is the Cartan structure equations. These are equations that view
geometric variables as differential forms and relates them to the
vielbein and spin connection through exterior derivatives and wedge products.

We first define the vielbein one form and the spin connection one form as:
\begin{align}
  e^a = e_\mu{}^a dx^\mu, \quad \omega^a{}_b = \omega_\mu{}^a{}_b dx^\mu
\end{align}
Then we can rewrite the tortion tensor and the Riemann curvature
tensor as differential forms as well:
\begin{align}
  T^a = \frac{1}{2} T^a{}_{\mu\nu} dx^\mu \wedge dx^\nu, \quad R^a{}_b =
  \frac{1}{2} R^a{}_{b\mu\nu} dx^\mu \wedge dx^\nu
\end{align}
where torsion tensor is defined as $ T^a{}_{\mu\nu} = 2
\Gamma^\lambda_{[\mu\nu]} e_\lambda{}^a $ and the Riemann curvature
tensor is defined as $ R^a{}_{b\mu\nu} = e^\lambda{}_b e_\rho{}^a
R^\rho{}_{\lambda\mu\nu} $.

The Cartan structure equations are:
\thm{
  Cartan's structure equation
  \begin{align}
    T^a = de^a + \omega^a{}_b \wedge e^b, \quad R^a{}_b = d\omega^a{}_b
    + \omega^a{}_c \wedge \omega^c{}_b
  \end{align}
}
The frist equation can be proved by directly using the relation
between the spin connection and the normal connection \cref{relationconnection}.
The second is a bit more complicated.

If we take another exterior derivative of both sides of the Cartan
structure equations, we can get the following identities:
\begin{align}
  dT^a + \omega^a{}_b \wedge T^b = R^a{}_b \wedge e^b, \quad dR^a{}_b +
  \omega^a{}_c \wedge R^c{}_b - R^a{}_c \wedge \omega^c{}_b = 0
\end{align}
The first one is called the Bianchi identity for torsion, and the
second one is just the normal Bianchi identity.

\subsection{Christoffel Symbol}

Above discussion we use general connection, now we consider the
special case of the Christoffel symbol, which is the unique torsion
free connection compatible with the metric.

We can see that:
\begin{itemize}
  \item Torsion free condition means that:
    \begin{align}
      de^a + \omega^a{}_b \wedge e^b = 0
    \end{align}
  \item Metric compatibility condition means that:
    \begin{align}
      \nabla g = 0 \quad \Rightarrow \quad \nabla_\mu \eta_{ab} = 0
    \end{align}
    Thus we can get that:
    \begin{align}
      \omega_{\mu ab} = -\omega_{\mu ba}
    \end{align}
\end{itemize}

\subsection{Extra: Convention and Definition of Differential Forms}

Here list out some definitions for differential forms, in case I make
a mistake in coefficients and signs.
\YL{[I'm not sure what convention Witten uses. Yet its not
    conceptually important and I use proportional to instead of equal to
in the main text ]}
